\section{Discussion}

Explicit token IDs allow separate tests of lexical recognition, ordered
transcription, and downstream correction.  The isolated-sign evaluation tests
whether the visual input suffices to recover one label.  The sequence benchmark
adds the need to recognize several labels, emit them in order, and stop at the
right point.  Its results support ordered transcription on controlled streams
from held-out signers.  The two evaluations use different checkpoints, so the
gap in accuracy cannot isolate the effect of sequence length.

The exploratory corruption audit gives no evidence of exact recovery at the
tested rates without sentence context.  Assessing correction for sentence
translation requires natural sentences with enough linguistic evidence to make
the intended word recoverable.

Discrete IDs trade flexibility for auditability.  Unlike continuous prompts, they
do not communicate a calibrated posterior over alternative tokens.  Future work
should compare top-$k$ or posterior-weighted interfaces, as explored in speech
prompting \citep{deng2025wav2prompt,ma2025legoslm}, while retaining a separately
scored visual transcript.  Such uncertainty may give the LLM useful alternatives
without returning to unconstrained embedding imitation.
